\subsectionaddtolist{آ}

برای اثبات این‌که مسئله‌ی \textsf{SUBSET-SUM} یک مسئله‌ی NP-complete است، باید دو بخش زیر را نشان دهیم:

{1. عضویت در کلاس \lr{NP}:}
ماشین تورینگ غیرقطعی $N$ را روی ورودی $\langle S, t \rangle$ که $S = \{x_1, \dots, x_k\}$ است به صورت زیر تعریف می‌کنیم:
\begin{enumerate}
    \item برای هر $i \in \{1, \dots, k\}$، به صورت غیرقطعی حدس بزن که آیا $x_i$ در زیرمجموعه انتخاب شود یا خیر.
    \item اعداد انتخاب‌شده را با هم جمع کن و مجموع را با $t$ مقایسه کن.
    \item اگر مجموع برابر $t$ بود، {بپذیر}؛ در غیر این صورت {رد کن}.
\end{enumerate}
ورودی پذیرفته می‌شود اگر و تنها اگر زیرمجموعه‌ای با مجموع $t$ وجود داشته باشد. حدس زدن زیرمجموعه و جمع کردن حداکثر $k$ عدد در زمان چندجمله‌ای برحسب طول ورودی انجام می‌شود؛ بنابراین $\text{\textsf{SUBSET-SUM}} \in \text{NP}$.

{2. اثبات \lr{NP-Hard} بودن (کاهش از \lr{3SAT}):}
مسئله‌ی \lr{3SAT} را به \lr{SUBSET-SUM} کاهش می‌دهیم ($\text{\lr{3SAT}} \le_P \text{\lr{SUBSET-SUM}}$).
فرض کنید فرمول فرم نرمال عطف‌دار $\phi$ دارای $n$ متغیر $x_1, \dots, x_n$ و $m$ بند $C_1, \dots, C_m$ باشد.
مجموعه‌ای از اعداد $S$ و هدف $t$ را در مبنای ۱۰ (یا مبنای ۴) با $n+m$ رقم به صورت زیر می‌سازیم:
\begin{itemize}
    \item برای هر متغیر $x_i$، دو عدد $v_i$ و $v'_i$ می‌سازیم. رقم $i$-ام آن‌ها برابر ۱ است. اگر $x_i$ در بند $C_j$ ظاهر شده باشد، رقم $(n+j)$-ام عدد $v_i$ برابر ۱ می‌شود، و اگر $\neg x_i$ در بند $C_j$ باشد، رقم $(n+j)$-ام $v'_i$ برابر ۱ می‌گردد.
    \item برای هر بند $C_j$، دو عدد کمکی $s_j$ و $s'_j$ می‌سازیم که تنها در رقم $(n+j)$-ام دارای مقدار ۱ هستند.
    \item عدد هدف $t$ به گونه‌ای تعریف می‌شود که $n$ رقم اول آن برابر ۱ (انتخاب دقیقا یکی از $x_i$ یا $\neg x_i$) و $m$ رقم بعدی آن برابر ۳ باشند.
\end{itemize}
این ساختار در زمان چندجمله‌ای برحسب طول فرمول ورودی ساخته می‌شود و $\phi$ صدق‌پذیر است اگر و تنها اگر زیرمجموعه‌ای وجود داشته باشد که مجموع آن برابر $t$ شود.
بنابراین \lr{SUBSET-SUM} مسئله‌ای \lr{NP-hard} و در نتیجه {\lr{NP-complete}} است.

\pagebreak

\subsectionaddtolist{ب}

در نمایش دودویی یا دهدهی، یک عدد طبیعی $N$ با طول $O(\log N)$ بیت نمایش داده می‌شود. در کاهش ارائه شده از \lr{3SAT} به \lr{SUBSET-SUM}، مقادیر اعداد ساخته‌شده تا مرتبه‌ی $10^{n+m}$ بزرگ می‌شدند، اما طول نمایش دودویی آن‌ها $O(n+m)$ بود (که چندجمله‌ای است).

اما در نمایش یکانی، عدد $N$ به صورت رشته‌ای از $N$ یک نمایش داده می‌شود؛ یعنی طول ورودی متناظر با خود مقدار عددی $N$ است. اگر همان کاهش را به کار ببریم، نوشتن عدد $10^{n+m}$ به صورت یکانی نیازمند $10^{n+m}$ کاراکتر است که نمایی برحسب طول ورودی اولیه ($n+m$) خواهد بود. بنابراین این کاهش دیگر زمان چندجمله‌ای نخواهد بود و اثبات سخت بودن شکست می‌خورد.

مسئله‌ی مجموع زیرمجموعه‌ها را می‌توان با استفاده از برنامه‌نویسی پویا حل کرد:
\begin{itemize}
    \item جدول $DP[i][w]$ را طوری تعریف می‌کنیم که نشان دهد آیا با استفاده از زیرمجموعه‌ای از $i$ عدد اول می‌توان به مجموع $w$ رسید یا خیر.
    \item پیچیدگی زمانی این الگوریتم $O(n \cdot t)$ است که $n = |S|$ تعداد اعداد و $t$ هدف است.
\end{itemize}

در حالت نمایش یکانی، طول ورودی مسئله برابر است با:
\[
L = \text{length}(\text{input}) \ge t + \sum_{i=1}^k x_i
\]
زیرا نمایش $t$ به صورت یکانی خود به تنهایی شامل $t$ علامت است. بنابراین $t \le L$ و $n \le L$.

در نتیجه، زمان اجرای الگوریتم برنامه‌نویسی پویا برابر است با:
\[
O(n \cdot t) \le O(L \cdot L) = O(L^2)
\]
چون زمان اجرای الگوریتم چندجمله‌ای برحسب طول ورودی یکانی است، مسئله در زمان چندجمله‌ای حل می‌شود و بنابراین $\text{\lr{UNARY-SSUM}} \in \text{P}$.
